% =====================================================================
%  CAE VIỆT NAM — TEMPLATE GIÁO TRÌNH TOÁN (LUYỆN THI CSCA)
%  Biên dịch bằng XeLaTeX (khuyến nghị) hoặc LuaLaTeX.
%  Cách dùng: sao chép file này, đổi \chuyende, \modulename, \modulepercent
%  và thêm nội dung vào các mục "NỘI DUNG CHUYÊN ĐỀ".
% =====================================================================
\documentclass[12pt, a4paper]{report}

% ---------- NGÔN NGỮ & PHÔNG CHỮ ----------
\usepackage{fontspec}
\usepackage{polyglossia}
\setmainlanguage{vietnamese}
\setotherlanguage{english}
\setmainfont{Latin Modern Roman}
\newfontfamily\cjkfont{SimSun}[Scale=0.95]   % dùng cho thuật ngữ Hán
\newcommand{\cjk}[1]{{\cjkfont #1}}

% ---------- TOÁN HỌC ----------
\usepackage{amsmath, amssymb, amsthm, mathtools}
\usepackage{unicode-math}

% ---------- BỐ CỤC TRANG ----------
\usepackage[a4paper, top=2.5cm, bottom=2.5cm, left=2.8cm, right=2.2cm]{geometry}
\usepackage{fancyhdr}
\usepackage{titlesec}
\usepackage{tocloft}
\usepackage{enumitem}
\usepackage{graphicx}
\usepackage{booktabs, array, longtable}
\usepackage{xcolor}
\usepackage[most]{tcolorbox}
\usepackage[colorlinks=true, linkcolor=cae_blue, citecolor=cae_blue, urlcolor=cae_blue]{hyperref}

% ---------- BẢNG MÀU THƯƠNG HIỆU ----------
\definecolor{cae_blue}{RGB}{0, 51, 102}
\definecolor{cae_gold}{RGB}{196, 149, 60}
\definecolor{box_bg}{RGB}{245, 247, 250}
\definecolor{box_line}{RGB}{0, 51, 102}

% ---------- TIÊU ĐỀ CHƯƠNG / MỤC ----------
\titleformat{\chapter}[display]
  {\normalfont\bfseries\color{cae_blue}}
  {\Large\MakeUppercase{\chaptertitlename} \thechapter}
  {6pt}
  {\Huge}
\titleformat{\section}
  {\normalfont\Large\bfseries\color{cae_blue}}{\thesection}{0.6em}{}
\titleformat{\subsection}
  {\normalfont\large\bfseries\color{cae_blue!80!black}}{\thesubsection}{0.6em}{}

% ---------- ĐẦU / CHÂN TRANG ----------
\pagestyle{fancy}
\setlength{\headheight}{15pt}
\fancyhf{}
\renewcommand{\headrulewidth}{0.4pt}
\renewcommand{\footrulewidth}{0.2pt}
\fancyhead[L]{\small\color{cae_blue}\leftmark}
\fancyhead[R]{\small\color{cae_blue}CAE VIỆT NAM — Luyện thi CSCA}
\fancyfoot[C]{\small\thepage}

% ---------- MÔI TRƯỜNG ĐỊNH LÝ / VÍ DỤ ----------
\newtcolorbox{dinhnghia}[1][]{
  colback=box_bg, colframe=box_line, boxrule=0.6pt, arc=1pt,
  fonttitle=\bfseries, title={Định nghĩa #1}, breakable}
\newtcolorbox{dinhly}[1][]{
  colback=box_bg, colframe=cae_gold, boxrule=0.6pt, arc=1pt,
  fonttitle=\bfseries, title={Định lý / Tính chất #1}, breakable}
\newtcolorbox{vidu}[1][]{
  colback=white, colframe=black!40, boxrule=0.4pt, arc=1pt,
  fonttitle=\bfseries, title={Ví dụ #1}, breakable}
\newtcolorbox{chuy}[1][]{
  colback=cae_gold!10, colframe=cae_gold!70!black, boxrule=0.4pt, arc=1pt,
  fonttitle=\bfseries, title={Chú ý #1}, breakable}
\newtcolorbox{baitap}[1][]{
  colback=white, colframe=cae_blue!50, boxrule=0.4pt, arc=1pt,
  fonttitle=\bfseries, title={Bài tập tự luyện #1}, breakable}

% ---------- MỤC LỤC ----------
\renewcommand{\cftchapfont}{\bfseries\color{cae_blue}}
\renewcommand{\cftsecfont}{\color{black}}

% =====================================================================
\begin{document}

% ---------------------- TRANG BÌA ----------------------
\begin{titlepage}
  \centering
  \vspace*{1.5cm}
  {\Large\bfseries\color{cae_blue} CAE VIỆT NAM}\\[0.3cm]
  {\small Tài liệu luyện thi CSCA — Môn Toán}\\[2.5cm]

  {\color{cae_gold}\rule{0.5\linewidth}{0.8pt}}\\[0.8cm]
  {\Huge\bfseries\color{cae_blue} CHUYÊN ĐỀ~\thechapter}\\[0.4cm]
  {\LARGE\bfseries TÊN CHUYÊN ĐỀ \cjk{（中文名）}}\\[0.6cm]
  {\large Module: \emph{Tên module đề thi} \quad—\quad Tỉ trọng: \emph{XX\%}}\\[0.8cm]
  {\color{cae_gold}\rule{0.5\linewidth}{0.8pt}}\\[2.5cm]

  {\normalsize Biên soạn: Hoàng Dương — CAE Việt Nam}\\
  {\small info@caevietnam.vn}\\[0.3cm]
  {\small \today}

  \vfill
  {\footnotesize © CAE Việt Nam. Tài liệu lưu hành nội bộ, không sao chép dưới mọi hình thức khi chưa được phép.}
\end{titlepage}

\tableofcontents
\clearpage

% =====================================================================
\chapter{Tên chuyên đề \cjk{（中文名）}}
\label{ch:mau}

\section{Kiến thức trọng tâm}
% NỘI DUNG CHUYÊN ĐỀ — thay bằng kiến thức thật

\begin{dinhnghia}
  Phát biểu định nghĩa. Ký hiệu toán học đặt trong dấu đô-la, ví dụ \(A \cup B\), \(A \cap B\).
\end{dinhnghia}

\begin{dinhly}
  Phát biểu định lý / tính chất, có thể kèm chứng minh ngắn hoặc chỉ dẫn chứng minh.
\end{dinhly}

\begin{chuy}
  Lưu ý, lỗi thường gặp, hoặc mẹo giải nhanh.
\end{chuy}

\section{Ví dụ minh họa}

\begin{vidu}
  \textbf{Đề bài.} Nội dung ví dụ.

  \textbf{Lời giải.} Trình bày lời giải từng bước, có căn chỉnh công thức bằng \texttt{align}.
  \[
    x^2 + 2x - 3 = 0 \iff (x-1)(x+3) = 0 \iff x = 1 \ \text{hoặc}\ x = -3.
  \]
\end{vidu}

\section{Bài tập tự luyện}

\begin{baitap}
  \begin{enumerate}[label=\textbf{Câu \arabic*.}, leftmargin=*, itemindent=0pt]
    \item Đề bài câu 1.
    \item Đề bài câu 2.
  \end{enumerate}
\end{baitap}

% ---- BẢNG PHÂN LOẠI CÂU HỎI THEO ĐỘ KHÓ (tùy chọn, đặt cuối chuyên đề) ----
\section*{Phân bố độ khó trong đề thi}
\addcontentsline{toc}{section}{Phân bố độ khó trong đề thi}
\begin{center}
\begin{tabular}{@{}lccc@{}}
  \toprule
  \textbf{Mức độ} & \textbf{Dễ} & \textbf{Trung bình} & \textbf{Khó} \\
  \midrule
  Số câu & XX & XX & XX \\
  \bottomrule
\end{tabular}
\end{center}

\end{document}
